\section{Additional Experiments}
\subsection{Confounds in $\hat{\beta}$}
\label{app:confounds}

\begin{table*}
\caption{We account for confounds when computing $\hat{\beta}$, and find that $\hat{\beta}$ remains positive throughout.}
\label{tab:beta_confounds}
\begin{tabular}{llcrr}
\toprule
& & $\hat\beta$ (s.e.) & Events & $N$ \\
\midrule
\emph{Baseline} & all events & +0.480$^{***}$ (0.117) & 106 & 848 \\
\addlinespace
\emph{Comparator choice} & best rival & +0.480$^{***}$ (0.117) & 106 & 848 \\
& second-best rival & +0.485$^{***}$ (0.120) & 106 & 848 \\
& top-3 average & +0.521$^{***}$ (0.117) & 106 & 848 \\
\addlinespace
\emph{Release year} & 2025 & +0.455$^{***}$ (0.114) & 60 & 480 \\
& 2026 & +0.849$^{***}$ (0.096) & 41 & 328 \\
\addlinespace
\emph{Sample composition} & events A & +0.500$^{***}$ (0.155) & 53 & 424 \\
& events B & +0.331$^{***}$ (0.105) & 53 & 424 \\
\addlinespace
\emph{Event-level heterogeneity} & release-event fixed effects & +0.949$^{***}$ (0.115) & 106 & 848 \\
\midrule
\emph{Regression to the mean} & measured spurious slope & +0.099 (0.038) & 20 splits & 116,257 \\
& correction-implied slope & +0.139 & -- & -- \\
\bottomrule
\end{tabular}
\end{table*}

We next account for other confounds potentially skewing $\hat{\beta}$. 
In Table~\ref{tab:beta_confounds}, we account for confounds including the selection of comparator (e.g., best rival, second-best rival, and top-3 average) and release year (2025 vs. 2026, with 2024 having only $5$ events).
For both of these, we find that $\hat{\beta}>0.45$, showing a similar catch-up effect. 
We additionally test whether the result is driven by a subset of release events by splitting events into two alternating subsamples and computing \(\hat{\beta}\) separately on each. 
We find that \(\hat{\beta}>0.33\) in both halves of the split.
To account for event-specific factors shared across categories, such as differences in the overall magnitude of improvement across releases, we include a release-event fixed effect \(\delta_e\).  
The resulting estimate remains positive and increases to \(\hat\beta=0.949\). 

We finally consider mechanical regression to the mean, since noise in the outgoing model's rating enters both the measured gap and subsequent improvement. 
Because the historical leaderboard provides only rating snapshots, we cannot re-estimate \(\hat{\beta}\) using independent ratings for each release event; instead, we use raw Arena battles to estimate the magnitude of this mechanical effect. 
We retrieve 116,257 battles, randomly split them in half, and fit Bradley--Terry ratings separately on each half to obtain two independent estimates, \(\hat{\theta}^{A}\) and \(\hat{\theta}^{B}\), of each model's rating.
We regress \(\hat{\theta}^{B}-\hat{\theta}^{A}\) on the gap computed from half \(A\) alone. 
Because the two estimates measure the same underlying model performance, any systematic positive relationship should arise from sampling variation in \(\hat\theta^A\), since the two estimates are constructed from independent random halves of the same set of battles.
These battles cover only April--July 2025 and three of our eight categories (plus the overall category), so we use this exercise to estimate the magnitude of the bias rather than to re-estimate \(\hat{\beta}\). 
Averaged over 20 random splits, the spurious slope is \(0.099\), and the errors-in-variables correction predicts $0.139$. 
As a result, measurement error appears too small to mechanically generate the observed catch-up relationship.

\subsection{Decomposing Capability Gaps}
\label{app:strong_competitors}

A larger capability gap can arise either because a developer's outgoing model performs poorly in a capability or because a rival performs well.
To distinguish between these factors, we regress incoming capability on outgoing capability and the rival frontier
separately:
\[
\theta^{\mathrm{new}}_{ikt} = \alpha_{ik}+\lambda_t +\rho\,\theta^{\mathrm{old}}_{ikt} +\delta\,\theta^{\mathrm{rival}}_{ikt} +\epsilon_{ikt},
\]
where $\alpha_{ik}$ are developer-category fixed effects and $\lambda_t$ are quarter fixed effects. 
Unlike our main specification, we use quarter rather than category-quarter fixed effects so that variation in the rival frontier remains identified.

We estimate $\hat\rho=0.363$ and $\hat\delta=-0.078$ (event-clustered SE $=0.116$ for $\hat\delta$). 
Conditional on the outgoing model's capability, we do not find evidence that a stronger rival frontier predicts stronger incoming capability.
Combined with our finding that $\hat{\beta}>0$, our takeaway is that developers improve more in capabilities where they lag, with a larger gap being consistent more strongly with a weaker developer rather than a stronger rival frontier. 
Proposition~\ref{prop:PrimitiveCatchup} provides one mechanism capable of generating the observed catch-up pattern, as diminishing marginal returns makes it plausible for catch-up to occur, although our observational analysis cannot identify the underlying mechanism.
Taking a step back, shared evaluation provides similar incentives to developers, while diminishing returns potentially allows weaker models to catch-up. 